% Generated draft: requires mathematical and source review.
\documentclass[11pt]{article}
\usepackage[margin=25mm]{geometry}
\usepackage{fontspec}
\setmainfont{Latin Modern Roman}
\usepackage{amsmath,amssymb,mathtools,unicode-math}
\setmathfont{Latin Modern Math}
\usepackage{enumitem,needspace,xurl,hyperref}
\setlength{\parindent}{0pt}
\setlength{\parskip}{5pt}
\setlength{\emergencystretch}{3em}
\newcommand{\N}{\mathbb{N}}
\newcommand{\Z}{\mathbb{Z}}
\newcommand{\Q}{\mathbb{Q}}
\newcommand{\R}{\mathbb{R}}
\newcommand{\C}{\mathbb{C}}
\newcommand{\F}{\mathbb{F}}
\newcommand{\A}{\mathbb{A}}
\newcommand{\PP}{\mathbb{P}}
\newcommand{\E}{\mathbb{E}}
\newcommand{\eps}{\varepsilon}
\newcommand{\dd}{\,\mathrm{d}}
\begin{document}
\section*{UnsolvedMath ID: 30005618}
\textbf{Source problem code:} catalogue entry
\subsection{Definitions and mathematical statement}
Interpret the unit square as the open domain \(D=(0,1)^2\subset\R^2\). A convex subdomain is understood in the usual sense as a nonempty open convex set \(\Omega\subseteq D\). For each such domain \(U\), define its Neumann eigenvalues, counted with multiplicity, by
\[
0=\lambda_0(U)<\lambda_1(U)\leq\lambda_2(U)\leq\cdots.
\]
Here \(\lambda\) is a Neumann eigenvalue if there exists a nonzero real-valued function \(u\in H^1(U)\) such that
\[
\int_U \nabla u\cdot\nabla v\,\dd x
=\lambda\int_U uv\,\dd x
\qquad\text{for every }v\in H^1(U).
\]
This is the weak formulation of \(-\Delta u=\lambda u\) in \(U\), with homogeneous Neumann boundary condition \(\partial_\nu u=0\) on \(\partial U\); it does not require a classical normal derivative at corners.

The supplied source asks for the inequality
\[
\forall\,\Omega\subseteq D\text{ nonempty, open, and convex},
\qquad \mu_2(\Omega)\geq\mu_2(D).
\]
The title identifies these as Neumann eigenvalues, but the supplied record does not define their indexing. Consequently, the intended precise assertion remains unresolved between the following conventions:
\begin{itemize}
\item If the zero eigenvalue is indexed by \(1\), then \(\mu_2(U)=\lambda_1(U)\), and the assertion is
\[
\forall\,\Omega\subseteq D\text{ nonempty, open, and convex},
\qquad \lambda_1(\Omega)\geq\lambda_1(D).
\]
\item If the zero eigenvalue is indexed by \(0\), then \(\mu_2(U)=\lambda_2(U)\), and the assertion is
\[
\forall\,\Omega\subseteq D\text{ nonempty, open, and convex},
\qquad \lambda_2(\Omega)\geq\lambda_2(D).
\]
\end{itemize}
The preceding spectral convention must be recovered from the source before selecting one formulation. The supplied record labels the problem open and the reconstructed statement \texttt{corrected\_verified}; these are supplied curation labels, not independent verification by this draft.
\subsection{Short English statement}
The proposed inequality says that every convex subdomain of the unit square has a specified Neumann eigenvalue at least as large as the corresponding eigenvalue of the square. The supplied excerpt leaves unresolved whether \(\mu_2\) means the first positive eigenvalue or the second positive eigenvalue, counted with multiplicity.
\subsection{Sources}
\begin{itemize}
\item Supplied local UnsolvedMath registry record: permanent ID \(30005618\), exact source problem code \texttt{catalogue entry}, title \emph{Second Neumann Eigenvalues of Convex Subdomains of a Square}, in the Oberwolfach Reports Open Problems collection. The record gives proposed year \(2023\) and supplies no named proposer for this problem.
\item Supplied source citation: \emph{Geometric Spectral Theory (2024)}, Oberwolfach Reports, DOI \texttt{10.4171/owr/2023/36}, pp.~2017--2106. Source URL and literature-source label ``original OWR report'': \url{https://doi.org/10.4171/owr/2023/36}. The supplied background describes Oberwolfach Reports as records of mathematical workshops and their open-problem discussions.
\item Supplied literature source: Cavallina--Funano--Henrot--Lemenant--Lucardesi--Sakaguchi, \emph{Neumann extremum problems}, \url{https://arxiv.org/abs/2312.13747}.
\item Supplied literature assessment, dated exactly \texttt{2026-08-22}: the record labels the problem open, reports known existence results and proofs for symmetric convex families, and states that the inequality for every convex subdomain of the unit square remains open. This is supplied dated curation derived from the OWR material, not independently verified literature assessment; its time-sensitive status and cited claims require independent verification.
\item Supplied statement-verification note: the statement was reconstructed and checked against the official Oberwolfach report. This draft preserves that reported verification history without claiming to have visited or checked the report.
\end{itemize}
\end{document}
